\chapter*{Introduction}\label{ch:intro}
\addcontentsline{toc}{chapter}{INTRODUCTION}

\begin{itemize}
    %\item Central dogma.
    %\item Human genome completion
    \item Incentive now to resolve the Expression Atlas transcriptomics and proteomics (but transcriptomics so much easier)
\end{itemize}

\begin{comment}
\Cref{fig:dogma} is the reconstruction by \citet{Crick:1958} of
what was conceived circa 1958 and is now known as
the core of the central dogma of molecular biology.
It presents the relationship between these three types of macromolecules:
one gene will be transcribed from \DNA\ to one \mRNA,
which will be in turn translated in one protein.
This theory still holds true
even though the truth is not as linear as first thought.
It also explains why the assumptions of strong correlations between the \mRNA\
and the protein levels of a gene have lasted (are still lasting) so long.

\begin{figure}%[!htbp]
    \includegraphics[scale=0.45]{introduction/dogma}\centering
    \caption[Central dogma of molecular biology proposed by F. Crick circa 1958]%
    {\label{fig:dogma}\textbf{Central dogma of molecular biology proposed by
    F. Crick circa 1958.} The solid lines present the proven modes of
    information transfers, the dashed ones are the ones that were postulated but
    have yet to be established.
    Nowadays, we know that the arrow from \gls{RNA} to \gls{DNA}
    should be a solid one, even though this happens on particular occasions.}
\end{figure}
\end{comment}

Over the past few years, a large number of gene expression atlases based on
\Rnaseq\ have been released for many different organisms\footnote{E.g.,
mouse~\mycite{Wu2009-lw,Ringwald2012}, pig~\mycite{Freeman2012},
sheep~\mycite{Clark2017-mw}, plants as maize~\mycite{Stelpflug2016-sm},
vigna~\mycite{Yao2016-se},pigeonpea~\mycite{Pazhamala2017-ig},
parasite, \eg\ \species{Schistosoma mansoni}.}.
Many focus on Human, as a whole,
\eg\ see~\mycite{Krupp2012,Jimenez-Lozano2012,Uhlen2015,GTEx2013}
or on on specific aspects,
\eg\ the organogenesis in the embryos~\mycite{Gerrard2016-zu}.

On the other hand, there is also a considerable sustained effort
to compile these studies in new ressources for the community;
either by relying on the initial expression levels released with the original
studies --- as for example
\hFoCi{TISSUES}{https://tissues.jensenlab.org/Search}{Santos2015-rj} or
\hFoCi{Harmonizome}{https://amp.pharm.mssm.edu/Harmonizome/}{Harmonizome}
--- or
by reprocessing the raw data as
\hFoCi{Expression Atlas}{https://www.ebi.ac.uk/gxa/home}{EBIgxa}.

While similar microarrays-based projects exist\footnote{E.g.\ Gene Expression
Atlas for Human Embryogenesis~\mycite{Yi2010-az}, Atlas of human primary
cells~\mycite{Mabbott2013-xf}, Gene atlas of mouse and human protein-encoding
transcriptomes (now hosted by \soft{BioGPS}) \mycite{Su2004-kc},
Allen Brain Atlas~\mycite{Hawrylycz2012-mx}.},
there is a burst in the use \Rnaseq\ as the main technology to establish
these references\footnote{Indeed, microarrays straightforwardly exhibit
a large drawback compared to \Rnaseq:
they allow to measure a finite number of already identified genes
while \Rnaseq\ has the capacity of discovering new ones.}.



The incentive is that compared to microarrays, there should be a lot less
batch effects and that these studies, in addition to explore the transcriptome
for new (undiscovered until now) transcripts, the relationships (coexpression,
coregulation between them) due to the fact that there are not these biases (as
microarrays) these studies can be used as reference for later studies.


Problem: while they don't present some of the batch effects of the microarrays,
there was no real studies \emph{at that time} on how good is the \Rnaseq.

While \emph{batch effects} were well-known for microarrays,
at the time I started this project
the subject was little known.
Since then several other studies have been published on the matter.
Whenever relevant, I present and discuss those results them relation to my owns.
The first piece towards this goal was done by the~\mycite{seqcmaqc}.



\begin{comment}
Normal tissues sampled from various people are phenotypically very similar,
despite individual hereditary and environmental variations.
At molecular level,
technical noise, batch effects and physical limitations
are also generating variations.
\end{comment}

Despite individual hereditary and environmental variations,
normal tissue samples are phenotypically very similar.
At the molecular level, in addition to the technical noise,
more variations are also expected.

Microarray studies have illustrated that integrating data
from different assays (even from the same platform and laboratory)
is challenging~\mycite{Walsh2015-nf},
which also explain the abundance of methods and tools~\mycite{Tseng2012-if}.
While meta-analyses or cross-normalisation of microarrays are still
used~\mycite{Afroz2016-xk} and they may be very pertinent and
useful,
microarray technology is subject to saturation effects and is semi-quantitative
while suffering from great batch effects~\mycite{Rung2013-ul}.

\Rnaseq\ nurtures the hope of enabling true quantitative measurements.


In \Paper{seqcmaqc}~\mycite{seqcmaqc} many facts that
we had observed within the lab have been reported.
Indeed, this paper highlights several consistent points
across sites and platforms,
such as:
\begin{itemize}[topsep=0pt,nolistsep,noitemsep]
        \setlength{\itemsep}{1pt}
        \setlength{\parskip}{0pt}
        \setlength{\parsep}{0pt}
    \item Detection of the same unannotated isoforms (provided a significant expression)
    \item Absolute quantification seems unachievable yet,
        but relative quantification is robust and even agrees well with \gls{qPCR}
        and microarrays.
    \item Comprehensive gene model crucial for accurate expression profiling.
    \item They advise to profile multiple libraries or samples per conditions.
    \item Gene-specific biases are observed for every platforms
        (including \gls{qPCR}).
    \item To achieve absolute quantification, it would need to estimate and
        remove background noise (improves accuracy but at the expense of precision)
    \item Pipelines have an impact on the final results.
\end{itemize}

Their results are sound as they have used the same samples across the different
sites and platforms which are very different biologically (Universal
Human Reference RNA) and Human Brain Reference RNA
(see~\mycite{MAQC_Consortium2006-hs}) (from the \gls{MAQC} consortium).

\section*{Aims of the thesis}

One main focus of my work for this thesis was to appraise the
consistency of findings that have been enlightened in normal Human tissues by
individual large scale transcriptomics and, more recently, large-scale
proteomics studies.

While the amount of data to process and integrate can be challenging,
paradoxically the data available to draw statistically significant conclusions
is often still too sparse. However, we do have enough to provide at least a
qualitative assessment of the consistency and evaluate the feasibility of a
quantitative assessment in general, if not, realise a quantitative integration
for a few key tissues.

\subsection*{Achievements}
\begin{itemize}
    \item First time so many datasets have been reprocessed together with the
        same annotation and hence comparison are better (y'a le truc de chez Burges,
        mais ils se sont concentré sur une petite partie des genes)
    \item surtout comparison of Gtex and Uhlen (aussi approfondi)
    \item reprocessing of all the proteome untargeted in human
    \item comparison entre protein expression and 2 datasets (meilleur que le
        science report dans le sens où il y a pas de polémique)
    \item siteweb//application de visualisaiton de la comparison
    \item liste croisé de tissues specifique et house keeping à travers datasets
       transcriptomic et protoemic
  \item attempt de normalisation basée sur les housekeepings
\end{itemize}


\begin{comment}
The core of central dogma of molecular biology --- \ie\ one \DNA\ (coding) gene
will be transcribed as one \mRNA\ transcript and will be in turn translated in
one protein --- still holds true even though we now know the truth is not as
linear as we thought first. Many regulation processes are implied before
and after the transcription and translation steps.
\Cref{fig:dogma} is the reconstruction by~\citet{Crick:1958} of
what was conceived circa 1958: the solid lines present the proven modes of
information transfers, the dashed ones are the ones that were postulated but
have yet to be established. To update this figure to our current knowledge, the line
starting from the RNA to the DNA should be a solid one.


\begin{figure}%[!htbp]
    \includegraphics[scale=0.6]{integration/dogma}\centering
    \caption[Central dogma of molecular biology proposed by F. Crick circa 1958]
    {\label{fig:dogma}\textbf{Central dogma of molecular biology proposed by
    F. Crick circa 1958.} The solid lines are proved direct transfer, the dashed
    ones the still have to be established. Nowadays, we know that the arrow from RNA
    to DNA should be a solid one, even though this happens on particular occasions.}
\end{figure}


However, as we still have to prove that proteins can be produced directly
from \DNA\footnote{which seems pretty unlikely to me},
it is assumed that translation creates each protein,
hence studying
transcriptomics and proteomics together for the same conditions should give us
insights on these regulatory processes.
Whereas it was implicitly assumed for a long time that for similar conditions
there should be a proportional relationship between transcriptomics and
proteomics, many studies\
\TK{add references, see TAC report 3}
done in cells have failed to show a high correlation between the two biological
layers.
From then, while there are still studies done jointly on transcriptomics and proteomics,
the focus has considerably shift either on the presence or absence of
a protein/transcript in a specific condition or if both transcript and protein
are differentially expressed between two conditions in the same way.
Nowadays, there are not many efforts to directly link their actual expression.
\end{comment}


